\documentclass[12pt]{book}
\usepackage[utf8]{inputenc}
\usepackage[british]{babel}
\usepackage[authoryear,datebegin]{flexbib}		%Uso de flexbib para la bibliografía.
\usepackage{graphics}
\usepackage{tikz}
\usepackage{float}
\usepackage{textcomp}
\usepackage{dsfont}
\usepackage{booktabs}		%Mejora el uso de las tablas (elimina las líneas verticales)
\usepackage{makeidx}		%Incluir Índice Analítico
%\usepackage[colorlinks=true, pdfborder={0 100 100}, linkcolor=orange, linktoc=page, citecolor=Green]{hyperref}
\usepackage{microtype}		%Arregla el aspecto gráfico del documento (incluye una mejor separación entre letras y palabras)
%\usepackage{amssymb}
%\usepackage{amsmath}
\usepackage{fancyhdr}
\usepackage[utf8]{inputenc} 

%Nuevos

\usepackage{rotating}
\usepackage{adjustbox}   
\usepackage{eurosym}
\usepackage{graphicx}
%\usepackage{caption}
%\usepackage{subcaption}
\usepackage{algorithm}
\usepackage{algpseudocode}
\usepackage{listings}

\usepackage{hyperref}
\usepackage{acronym}

\usepackage{pgf,tikz}
\usetikzlibrary{arrows}
\pagestyle{empty}
\newcommand{\degre}{\ensuremath{^\circ}}


\usepackage[compact]{titlesec}
\titlespacing{\section}{2pt}{*2}{*2}
\titlespacing{\subsection}{2pt}{*2}{*2}
\titlespacing{\subsubsection}{2pt}{*2}{*2}
\titlespacing{\paragraph}{2pt}{*2}{*2}
\titlespacing{\subparagraph}{2pt}{*2}{*2}

%\setlength{\intextsep}{20pt plus 2.0pt minus 2.0pt}

%\setlength{\textfloatsep}{1pt}

\definecolor{orange}{RGB}{250,120,0}
\definecolor{Green}{RGB}{50,200,100}

\interfootnotelinepenalty=10000

\setlength{\parskip}{0.7mm}

\newcommand{\N}{\mathds{N}}	%Redefinición símbolo números Naturales
\newcommand{\Z}{\mathds{Z}}	%Redefinición símbolo números Enteros
\newcommand{\Q}{\mathds{Q}}	%Redefinición símbolo números Racionales
\newcommand{\I}{\mathds{I}}	%Redefinición símbolo números Irracionales
\newcommand{\R}{\mathds{R}}	%Redefinición símbolo números Reales
\newcommand{\p}{\ldotp}		%Redefinición símbolo 
\newtheorem{theorem}{Theorem}
\newtheorem{teorema}{Teorema}
\newtheorem{lema}{Lema}
\newtheorem{definicion}{Definición}
\newtheorem{ax}{Axioma}
\newtheorem{corolario}{Corolario}
\renewcommand{\rmdefault}{phv} % Arial
\renewcommand{\sfdefault}{phv} % Arial
\renewcommand{\baselinestretch}{1.2}
\renewcommand{\lstlistingname}{Listado}

\makeindex

\textheight=21.5cm    %alto del texto
\textwidth=15cm     %ancho del texto	
\topmargin=0cm     %margen superior
\oddsidemargin=1.5cm  %margen izquierdo (por defecto queda a 1.5 cm)

\usepackage{multirow}
\setcounter{tocdepth}{3} 	%Inserta subsubsecciones y párrafos en el indice
\usepackage{cleveref}		%Inserta la palabra correspondiente al referenciar un objeto (Figura, Tabla, etc)
\setlength{\parindent}{12pt}
\crefname{equation}{ecuación}{ecuaciones}
\crefname{table}{tabla}{tablas}
\crefname{figure}{figura}{figuras}
\crefname{section}{sección}{Sección}
\crefname{chapter}{capítulo}{Capítulo}

\graphicspath{{Imagenes/}} 	%La carpeta de dibujos es figs
%\includeonly{Capitulo_04}	%Incluye sólo los archivos que se especifican como parámetro. Se mantiene numeración global del documento.

% encabezados
\lhead[\thepage]{CHAPTER \thechapter}
\chead[]{}
\rhead[CHAPTER \thechapter]{\thepage}
\renewcommand{\headrulewidth}{0.5pt}

% pie de pagina
\lfoot[]{}
\cfoot[]{}
\rfoot[]{}
\renewcommand{\footrulewidth}{0pt}

% primera pagina de un capitulo
\fancypagestyle{plain}{
	\fancyhead[L]{}
	\fancyhead[C]{}
	\fancyhead[R]{\thepage}
	\fancyfoot[L]{}
	\fancyfoot[C]{}
	\fancyfoot[R]{}
	\renewcommand{\headrulewidth}{0pt}
	\renewcommand{\footrulewidth}{0pt}
}

%\pagestyle{fancy}
%
\lstset{
	language=XML,    
	basicstyle=\scriptsize,
	upquote=true,
	aboveskip={1.0\baselineskip},
	columns=fullflexible,
	showstringspaces=false,
	extendedchars=true,
	breaklines=true,
	showtabs=false,
	showspaces=false,
	showstringspaces=false,
	identifierstyle=\ttfamily,
	keywordstyle=\color[rgb]{0,0,1},
	commentstyle=\color[rgb]{0.133,0.545,0.133},
	stringstyle=\color[rgb]{0.627,0.126,0.941},
	frame=single,
}

\lstdefinelanguage{XML2}
{
basicstyle=\scriptsize,
upquote=true,
aboveskip={1.0\baselineskip},
columns=fullflexible,
showstringspaces=false,
extendedchars=true,
breaklines=true,
showtabs=false,
showspaces=false,
showstringspaces=false,
identifierstyle=\ttfamily,
keywordstyle=\color[rgb]{0,0,1},
commentstyle=\color[rgb]{0.133,0.545,0.133},
stringstyle=\color[rgb]{0.627,0.126,0.941},
frame=single
}




\begin{document}
	
\renewcommand{\contentsname}{Index}	%Cambiar nombre de Índice por Tabla de contenidos
\renewcommand{\listfigurename}{Index of figures}	%Cambiar nombre de Índice de Figuras por Índice de Ilustraciones
\renewcommand{\indexname}{Índice analítico}		%Cambiar nombre de Índice alfabético -> Índice analítico
% \renewcommand{\figurename}{Índice de Ilustraciones}	%Cambiar nombre de Índice de Figuras por Índice de Ilustraciones
\renewcommand{\listtablename}{Index of tables}
\renewcommand{\tablename}{Table}

%%%%%%%%%%%%%%%%%%%%%%%%%%%%%%%%%%
    %PÁGINAS PRELIMINARES
%%%%%%%%%%%%%%%%%%%%%%%%%%%%%%%%%%
%(1) PORTADA 			[L]
%(2) DERECHOS DE AUTOR 		[L]
%(3) HOJA DE CALIFICACIONES	[]
%(4) DEDICATORIA(OPTATIVA)	[]
%(5) AGRADECIMIENTOS(OPTATIVA)	[]
%(6) TABLA DE CONTENIDOS	[]
%(7) ÍNDICE DE TABLAS		[]
%(8) ÍNDICE DE ILUSTRACIONES	[]
%(9) RESUMEN			[]
%%%%%%%%%%%%%%%%%%%%%%%%%%%%%%%%%%

%%%%%%%%%%%%%%%%%%%%%%%%%%%%%%%%%%
%(1) PORTADA			[L]
%%%%%%%%%%%%%%%%%%%%%%%%%%%%%%%%%%
\begin{titlepage}
    
\begin{center}    
    \includegraphics[width=10.0cm]{images/initUJI.png}
\end{center}


\vspace*{2cm}

  \begin{center}
    \Large{\textbf{Thesis title}}\\ 
  \end{center}
  
\vspace*{0.5cm}
  
  \begin{center}
    Doctoral Thesis\\ 
    \textbf{Your name}
  \end{center}

\vspace*{1cm}

  \begin{center}
    \begin{tabular}{ll}
      \textbf{Supervisors:}& Dr. Supervisor 1 (University 1)\\
      & Dr. Supervisor 2 (University 2)\\
      & Dr. Supervisor 3 (University 3)\\
    \end{tabular}
    \end{center}
\vspace{3.5cm}

  \begin{center}
  Dissertation submitted to the Institute of New Imaging Technologies (INIT) in partial fulfilment of the requirements for the Degree of Doctor by the University Jaume I
\end{center}
  %\end{flushright}

\vspace{1.5cm}

  \begin{center}
    \textbf{Castelló de la Plana (Spain)\\november 2018}
  \end{center}
\end{titlepage}

\newpage
\mbox{}
\thispagestyle{empty} % para que no se numere esta página
\newpage
\thispagestyle{empty}

%%%%%%%%%%%%%%%%%%%%%%%%%%%%%%%%%%
%(2) DERECHOS DE AUTOR		[L]
%%%%%%%%%%%%%%%%%%%%%%%%%%%%%%%%%%
\vspace*{18cm}

\begin{tabular}{ll}
		\multirow{5}{*}{\includegraphics[width=2.7cm]{images/eu.png}} & \\
		  &This dissertation is funded by the European Commission \\
		  &within the Marie Skłodowska-Curie Actions (ITN-EJD). Grant\\
          &Agreement num. 642332 - GEO-C - H2020-MSCA-ITN-2014. 
\end{tabular}


%%%%%%%%%%%%%%%%%%%%%%%%%%%%%%%%%%
%(3) HOJA DE CALIFICACIONES	[]
%%%%%%%%%%%%%%%%%%%%%%%%%%%%%%%%%%

%%%%%%%%%%%%%%%%%%%%%%%%%%%%%%%%%%
%(4) DEDICATORIA(OPTATIVA)	[]
%%%%%%%%%%%%%%%%%%%%%%%%%%%%%%%%%%

\newpage
\thispagestyle{empty} 
\mbox{}
\newpage
\thispagestyle{empty} % para que no se numere esta página
\pagenumbering{roman}% para comenzar la numeración de paginas en números romanos
\thispagestyle{empty}


\chapter*{}
\begin{flushright}
\textit{My family}
\end{flushright}


\newpage
\mbox{}
%\thispagestyle{empty} % para que no se numere esta página



%%%%%%%%%%%%%%%%%%%%%%%%%%%%%%%%%%
%(5) AGRADECIMIENTOS(OPTATIVA)	[]
%%%%%%%%%%%%%%%%%%%%%%%%%%%%%%%%%%
\chapter*{Acknowledgments} % *si no queremos que añada la palabra "Capitulo"
% \addcontentsline{toc}{chapter}{Agradecimientos} % si queremos que aparezca en el índice

bla bla



%%%%%%%%%%%%%%%%%%%%%%%%%%%%%%%%%%

%%%%%%%%%%%%%%%%%%%%%%%%%%%%%%%%%%
%(9) RESUMEN			[]
%%%%%%%%%%%%%%%%%%%%%%%%%%%%%%%%%%
\newpage
\chapter*{Resumen}
\addcontentsline{toc}{chapter}{Resumen} % si queremos que aparezca en el índice

Spanish thesis summary (all thesis with 5 pages)

\vspace{1.0cm}

\textbf{Palabras clave}: \textit{blabla1}, \textit{blabla2}, \textit{blabla3}. 


\newpage
\chapter*{Abstract}
\addcontentsline{toc}{chapter}{Abstract} % si queremos que aparezca en el índice

bla bla

\vspace{1.0cm}

\textbf{Keywords}: \textit{blabla1}, \textit{blabla2}, \textit{blabla3}. 


%\textbf{Keywords}: \textit{Sensores}, \textit{Internet of Things}, \textit{Publicación automatizada}, \textit{Interoperabilidad}, \textit{Estándares GIS}, \textit{Análisis en Tiempo Real}. 

%%%%%%%%%%%%%%%%%%%%%%%%%%%%%%%%%%
%(6) TABLA DE CONTENIDOS 	[]
%%%%%%%%%%%%%%%%%%%%%%%%%%%%%%%%%%
\newpage
\tableofcontents
\newpage
%\newrgbcolor{xdxdff}{0.49 0.49 1}
%\newrgbcolor{ffqqtt}{1 0 0.2}
%\psset{xunit=1.0cm,yunit=1.0cm,algebraic=true,dotstyle=*,dotsize=3pt
%0,linewidth=0.8pt,arrowsize=3pt 2,arrowinset=0.25}

%%%%%%%%%%%%%%%%%%%%%%%%%%%%%%%%%%
%(7) ÍNDICE DE TABLAS		[]
%%%%%%%%%%%%%%%%%%%%%%%%%%%%%%%%%%

%%%%%%%%%%%%%%%%%%%%%%%%%%%%%%%%%%
%(8) ÍNDICE DE ILUSTRACIONES	[]
\newpage
\listoffigures

\newpage
\chapter*{Acronyms}
\addcontentsline{toc}{chapter}{Acronyms} 
\markboth{Acronyms}{Acronyms}
\begin{acronym}[AWGN]
	\acro{IP}{Protocolo de Internet}
	\acro{IoT}{Internet de las Cosas}
	\acro{URI}{Identificador de Recursos Uniforme} 
	\acro{WSN}{Redes de Sensores Inalámbricos} 
	\acro{GIS}{Sistemas de Información Geoespacial}
\end{acronym}

\newpage
\pagenumbering{arabic}

\include{chapter01}		
\include{conclusions}	


\include{Anexo}



%(2) BIBLIOGRAPHY			[]
\newpage
\bibliography{biblio}
\bibliographystyle{flexbib}
% \nocite{KasnerNewman, Martinez, Niven, Kaplun, Socioepistemologia, Hitt2007} Tener referencia en la bibliografía que no ha sido citada.

\end{document}